\chapter{Interchanging Limits, Derivatives and Integrals}
\label{ch:ii12}

Passing a limit through a derivative or integral requires more than pointwise convergence. The correct hypotheses combine uniform control with convergence of derivative data or domination of integral errors.

\section*{Learning goals}
After this chapter, the reader should be able to:
\begin{itemize}
\item verify hypotheses that justify interchanging a limit with an integral or derivative;
\item prove convergence of derivatives by combining uniform convergence of derivatives with convergence at one base point;
\item construct examples showing that pointwise convergence alone does not justify differentiating or integrating termwise;
\item estimate remainders uniformly enough to pass a limit through an operation;
\item distinguish sufficient hypotheses for integral interchange from those needed for derivative interchange;
\item use a theorem only after checking its domain, regularity, and convergence assumptions explicitly;
\end{itemize}
\section*{Conceptual roadmap}
\[
\boxed{\text{definition}\;\longrightarrow\;\text{estimate}\;\longrightarrow\;\text{theorem}\;\longrightarrow\;\text{application}.}
\]

\section{Limits and integrals}
Uniform convergence allows Riemann integration to commute with limits on compact intervals.


% BEGIN VOL02-EXPANSION II12-example-01
\begin{example}[Uniform convergence justifies integral interchange quantitatively]\label{ex:ii12-expand-01}
Suppose \(f_n\to f\) uniformly on \([a,b]\). Then
\[
\left|\int_a^b f_n-\int_a^b f\right|
\le \int_a^b|f_n-f|
\le (b-a)\|f_n-f\|_\infty.
\]
If the sup error is at most \(10^{-4}\) on an interval of length \(7\), then
the integral error is at most \(7\times10^{-4}\). Thus the interchange theorem
comes with a direct stability estimate, not merely a qualitative limit.
\end{example}
% END VOL02-EXPANSION II12-example-01
\section{Integral error estimates}
The difference of integrals is controlled by interval length times the sup-norm error.

\section{Limits and derivatives}
Uniform convergence of derivatives plus convergence at one base point yields a differentiable limit.


% BEGIN VOL02-EXPANSION II12-example-02
\begin{example}[Derivative convergence from one anchor value]\label{ex:ii12-expand-02}
Let
\[
f_n(x)=c_n+\int_0^x g_n(t)\,dt
\]
on \([0,1]\), where \(c_n\to c\) and \(g_n\to g\) uniformly. Then
\[
|f_n(x)-f_m(x)|
\le |c_n-c_m|+\|g_n-g_m\|_\infty
\]
uniformly in \(x\). Hence \(f_n\) converges uniformly to
\[
f(x)=c+\int_0^x g(t)\,dt,
\]
and \(f'=g\). The single anchor value fixes the additive constant that
derivative information alone cannot determine.
\end{example}
% END VOL02-EXPANSION II12-example-02
\section{Why derivative exchange is harder}
Differentiation amplifies small-scale oscillation and therefore needs stronger hypotheses.

\section{Series versions}
The same principles apply to term-by-term integration and differentiation of function series.

\section{Uniform convergence on compacta}
Local uniform convergence is often the right compromise on noncompact domains.

\section{Counterexamples}
Pointwise convergence alone can fail spectacularly for derivatives and integrals.


% BEGIN VOL02-EXPANSION II12-example-03
\begin{example}[A narrow spike defeats pointwise integral interchange]\label{ex:ii12-expand-03}
Define a triangular continuous spike \(f_n\) on \([0,1]\) supported in
\([0,2/n]\), with height \(n\) at \(1/n\). Its area is
\[
\frac12\cdot\frac{2}{n}\cdot n=1.
\]
For every fixed \(x>0\), eventually \(x\) lies outside the support, while
\(f_n(0)=0\). Thus \(f_n(x)\to0\) pointwise everywhere, but
\[
\int_0^1 f_n=1
\]
for all \(n\). Pointwise convergence alone gives no uniform control over the
moving concentration.
\end{example}
% END VOL02-EXPANSION II12-example-03
\section{Stability viewpoint}
Interchange theorems are continuity statements for integral or derivative operators in suitable function-space topologies.

\section{Core structural results}

\begin{theorem}[Uniform limit under the integral]\label{thm:ii12-01}
If $f_n\to f$ uniformly on $[a,b]$ and each $f_n$ is Riemann integrable, then $f$ is integrable and $\int f_n\to\int f$.
\begin{proof}
The integral difference is at most $(b-a)\|f_n-f\|_\infty$, which tends to zero; integrability follows by approximating $f$ uniformly with an integrable $f_n$.
\end{proof}
\end{theorem}

\begin{theorem}[Derivative limit theorem]\label{thm:ii12-02}
If $f_n\in C^1[a,b]$, $f_n(x_0)$ converges at one point, and $f_n'$ converges uniformly to $g$, then $f_n$ converges uniformly to a differentiable $f$ with $f'=g$.
\begin{proof}
Write $f_n(x)=f_n(x_0)+\int_{x_0}^x f_n'$. Uniform convergence of derivatives gives uniform Cauchy control of the integrals, and the fundamental theorem identifies the derivative of the limit.
\end{proof}
\end{theorem}

\begin{theorem}[Termwise integration]\label{thm:ii12-03}
If $\sum f_n$ converges uniformly and each $f_n$ is integrable, then the sum is integrable and $\int\sum f_n=\sum\int f_n$.
\begin{proof}
Apply the uniform-limit integral theorem to the partial sums.
\end{proof}
\end{theorem}

\section{Worked examples}

\begin{example}[Integral interchange estimate]\label{ex:ii12-worked-01}
Prove $|\int_a^b f_n-\int_a^b f|\le(b-a)\|f_n-f\|_\infty$. Apply $|\int g|\le\int|g|$ and bound the integrand by its supremum.
\end{example}

\begin{example}[A pointwise failure for integrals]\label{ex:ii12-worked-02}
Construct continuous spikes $f_n$ on $[0,1]$ converging pointwise to $0$ but with integrals not tending to $0$. Use triangular spikes of height $n$ and width $2/n$ centered near zero so area stays $1$ while each fixed positive point eventually lies outside the support.
\end{example}

\begin{example}[Termwise integration of a geometric series]\label{ex:ii12-worked-03}
For $|x|\le r<1$, integrate $\sum_{n\ge0}x^n$ termwise. The series is uniformly convergent by the M-test, so $\int_0^r\sum x^n dx=\sum r^{n+1}/(n+1)=-\log(1-r)$.
\end{example}

\begin{example}[Derivative exchange mechanism]\label{ex:ii12-worked-04}
Why is convergence at one base point needed? Derivatives determine functions only up to additive constants. Convergence at one point fixes those constants.
\end{example}

\section{Solved dossiers}
The dossiers are longer solved problems. Legacy mappings guide topic selection where explicit retained source blocks exist; otherwise fresh canonical problems complete the chapter.

\begin{problem}[Integral interchange estimate]\label{prob:ii12-dossier-01}
Prove $|\int_a^b f_n-\int_a^b f|\le(b-a)\|f_n-f\|_\infty$.
\end{problem}
\begin{solution}
Apply $|\int g|\le\int|g|$ and bound the integrand by its supremum.
\end{solution}

\begin{problem}[A pointwise failure for integrals]\label{prob:ii12-dossier-02}
Construct continuous spikes $f_n$ on $[0,1]$ converging pointwise to $0$ but with integrals not tending to $0$.
\end{problem}
\begin{solution}
Use triangular spikes of height $n$ and width $2/n$ centered near zero so area stays $1$ while each fixed positive point eventually lies outside the support.
\end{solution}

\begin{problem}[Termwise integration of a geometric series]\label{prob:ii12-dossier-03}
For $|x|\le r<1$, integrate $\sum_{n\ge0}x^n$ termwise.
\end{problem}
\begin{solution}
The series is uniformly convergent by the M-test, so $\int_0^r\sum x^n dx=\sum r^{n+1}/(n+1)=-\log(1-r)$.
\end{solution}

\begin{problem}[Derivative exchange mechanism]\label{prob:ii12-dossier-04}
Why is convergence at one base point needed?
\end{problem}
\begin{solution}
Derivatives determine functions only up to additive constants. Convergence at one point fixes those constants.
\end{solution}

\begin{problem}[Derivative counterexample]\label{prob:ii12-dossier-05}
Let $f_n(x)=\sin(nx)/n$. What happens to $f_n$ and $f_n'$?
\end{problem}
\begin{solution}
$f_n\to0$ uniformly, but $f_n'=\cos(nx)$ does not converge pointwise at most points, showing uniform convergence of functions alone does not control derivatives.
\end{solution}

\begin{problem}[Uniform derivative convergence]\label{prob:ii12-dossier-06}
If $\|f_n'-g\|_\infty\to0$, estimate the derivative-integral error.
\end{problem}
\begin{solution}
For any $x$, $|\int_{x_0}^x(f_n'-g)|\le |b-a|\|f_n'-g\|_\infty$.
\end{solution}

\begin{problem}[Series of derivatives]\label{prob:ii12-dossier-07}
Suppose $\sum f_n'$ converges uniformly and $\sum f_n(x_0)$ converges. What follows?
\end{problem}
\begin{solution}
The series $\sum f_n$ converges uniformly to a differentiable function whose derivative is $\sum f_n'$.
\end{solution}

\begin{problem}[Local uniform convergence]\label{prob:ii12-dossier-08}
Why is uniform convergence on every compact interval enough for local integral interchange?
\end{problem}
\begin{solution}
Any fixed finite integration interval is compact, so the uniform theorem applies there.
\end{solution}

\begin{problem}[Continuity of integration operator]\label{prob:ii12-dossier-09}
View $I(f)=\int_a^bf$. What is its operator norm on $C[a,b]$ with sup norm?
\end{problem}
\begin{solution}
$|I(f)|\le(b-a)\|f\|_\infty$, with equality for the constant function $1$, so the norm is $b-a$.
\end{solution}

\begin{problem}[Differentiation is not bounded in sup norm]\label{prob:ii12-dossier-10}
Use $f_n(x)=\sin(nx)/n$ to explain.
\end{problem}
\begin{solution}
$\|f_n\|_\infty=1/n\to0$ while $\|f_n'\|_\infty=1$, so no fixed bound $\|f'\|_\infty\le C\|f\|_\infty$ can hold.
\end{solution}

\begin{problem}[M-test bridge]\label{prob:ii12-dossier-11}
Why does the Weierstrass M-test often precede termwise integration?
\end{problem}
\begin{solution}
It supplies uniform convergence of the function series, exactly the hypothesis needed to exchange sum and integral.
\end{solution}

\begin{problem}[Interchange synthesis]\label{prob:ii12-dossier-12}
What unifies limit/integral and limit/derivative theorems?
\end{problem}
\begin{solution}
Both ask whether an operator is continuous under a chosen function-space topology. Integration is continuous in sup norm; differentiation requires a stronger topology controlling derivatives.
\end{solution}

\section{Exercises with complete solutions}

\begin{exercise}\label{exr:ii12-01}
If $f_n\to f$ uniformly, do integrals converge?
\end{exercise}
\begin{hint}
On a compact interval, yes. For interchanging limit and integral, identify a theorem whose hypotheses give uniform or dominated control before passing the limit.
\end{hint}
\begin{solution}
Yes, for Riemann-integrable terms.
\end{solution}

\begin{exercise}\label{exr:ii12-02}
Bound the integral error by sup error.
\end{exercise}
\begin{hint}
Use interval length. For derivative interchange, check uniform convergence of the derivative sequence and convergence at one base point.
\end{hint}
\begin{solution}
At most $(b-a)\|f_n-f\|_\infty$.
\end{solution}

\begin{exercise}\label{exr:ii12-03}
Does uniform convergence imply derivative convergence?
\end{exercise}
\begin{hint}
Use $\sin(nx)/n$. Integrate the derivative convergence from a fixed base point to reconstruct convergence of the functions.
\end{hint}
\begin{solution}
No.
\end{solution}

\begin{exercise}\label{exr:ii12-04}
Why one base point in derivative theorem?
\end{exercise}
\begin{hint}
Constants. To build a counterexample, look for pointwise convergence with spikes or rapidly changing slopes that defeat uniform control.
\end{hint}
\begin{solution}
Derivative data omit additive constants.
\end{solution}

\begin{exercise}\label{exr:ii12-05}
Can a uniformly convergent series be integrated termwise?
\end{exercise}
\begin{hint}
Under integrability. Estimate the remainder uniformly in the variable before moving a limit through an operation.
\end{hint}
\begin{solution}
Yes.
\end{solution}

\begin{exercise}\label{exr:ii12-06}
What topology makes integration continuous?
\end{exercise}
\begin{hint}
Sup norm on a compact interval. Termwise differentiation is stronger than termwise integration; check the derivative hypotheses separately.
\end{hint}
\begin{solution}
The sup norm.
\end{solution}

\begin{exercise}\label{exr:ii12-07}
Is differentiation bounded in the sup norm?
\end{exercise}
\begin{hint}
Use high-frequency functions. Write the theorem's hypotheses next to the sequence at hand and verify each one explicitly.
\end{hint}
\begin{solution}
No.
\end{solution}

\begin{exercise}\label{exr:ii12-08}
If derivatives converge uniformly to $g$, what candidate derivative has the limit?
\end{exercise}
\begin{hint}
Use the theorem. If only pointwise convergence is known, avoid exchanging operations until an additional bound or compactness argument supplies control.
\end{hint}
\begin{solution}
$g$.
\end{solution}


% BEGIN VOL02-EXPANSION II12-exercises-01
\section{Graded supplementary exercises}
These exercises extend the preserved foundational set and are graded by mathematical role.

\subsection*{Standard computations}

\begin{exercise}[Integral error]\label{exr:ii12-expand-01}
If \(\|f_n-f\|_\infty\le1/n\) on \([0,3]\), bound the integral error.
\end{exercise}
\begin{hint}
Multiply the sup error by interval length.
\end{hint}
\begin{solution}
At most \(3/n\).
\end{solution}

\begin{exercise}[Base-point constants]\label{exr:ii12-expand-02}
If \(f_n'=0\) for all \(n\) but \(f_n(0)=n\), do the functions converge?
\end{exercise}
\begin{hint}
Derivatives determine functions only up to constants.
\end{hint}
\begin{solution}
No; \(f_n\equiv n\) diverges despite identical derivatives.
\end{solution}

\begin{exercise}[Termwise integral]\label{exr:ii12-expand-03}
For \(|x|\le r<1\), justify integrating \(\sum x^n\) termwise.
\end{exercise}
\begin{hint}
Use the M-test with \(r^n\).
\end{hint}
\begin{solution}
The series is uniformly convergent, so \(\int_0^r\sum x^n dx=\sum r^{n+1}/(n+1)\).
\end{solution}

\begin{exercise}[Uniform derivative limit]\label{exr:ii12-expand-04}
If \(f_n'(x)=x/n\) on \([0,1]\), what is the uniform limit of derivatives?
\end{exercise}
\begin{hint}
Take the sup norm.
\end{hint}
\begin{solution}
\(\|f_n'\|_\infty=1/n\to0\), so the limit is \(0\).
\end{solution}

\begin{exercise}[Local uniformity]\label{exr:ii12-expand-05}
Does \(1/(1+x/n)\to1\) uniformly on every compact interval?
\end{exercise}
\begin{hint}
Bound \(x\) on a fixed compact set.
\end{hint}
\begin{solution}
Yes; on \([-R,R]\) for large \(n\), the deviation is bounded by a constant times \(R/n\).
\end{solution}

\subsection*{Proofs}

\begin{exercise}[Integral interchange]\label{exr:ii12-expand-06}
Prove uniform convergence permits passing a limit through the Riemann integral.
\end{exercise}
\begin{hint}
Use the sup-norm integral estimate.
\end{hint}
\begin{solution}
The integral differences are bounded by \((b-a)\|f_n-f\|_\infty\to0\).
\end{solution}

\begin{exercise}[Derivative theorem anchor]\label{exr:ii12-expand-07}
Explain why convergence of \(f_n(x_0)\) at one point is needed in the derivative-limit theorem.
\end{exercise}
\begin{hint}
All derivatives are unchanged by additive constants.
\end{hint}
\begin{solution}
Without anchoring constants, functions can drift vertically and fail to converge even when derivatives converge perfectly.
\end{solution}

\begin{exercise}[Termwise integration of series]\label{exr:ii12-expand-08}
Prove a uniformly convergent series of integrable functions can be integrated termwise.
\end{exercise}
\begin{hint}
Apply the uniform integral theorem to partial sums.
\end{hint}
\begin{solution}
If \(S_N\to S\) uniformly, then \(\int S_N\to\int S\), while \(\int S_N=\sum_{n\le N}\int f_n\).
\end{solution}

\begin{exercise}[Uniform derivative Cauchy estimate]\label{exr:ii12-expand-09}
If \(f_n(x_0)\) is Cauchy and \(f_n'\) is uniformly Cauchy on \([a,b]\), prove \(f_n\) is uniformly Cauchy.
\end{exercise}
\begin{hint}
Write differences using the fundamental theorem from \(x_0\).
\end{hint}
\begin{solution}
\(|f_n(x)-f_m(x)|\le|f_n(x_0)-f_m(x_0)|+(b-a)\|f_n'-f_m'\|_\infty\).
\end{solution}

\subsection*{Counterexamples and hypothesis testing}

\begin{exercise}[Pointwise integral failure]\label{exr:ii12-expand-10}
Give a pointwise-zero sequence of continuous functions whose integrals stay \(1\).
\end{exercise}
\begin{hint}
Use narrowing triangular spikes with increasing height.
\end{hint}
\begin{solution}
Height \(n\), base \(2/n\) triangular spikes have area \(1\) and converge pointwise to zero if placed at the endpoint appropriately.
\end{solution}

\begin{exercise}[Derivative interchange failure]\label{exr:ii12-expand-11}
Why can \(f_n(x)=\sin(nx)/n\) converge uniformly while derivatives fail to converge?
\end{exercise}
\begin{hint}
Compute \(f_n'\).
\end{hint}
\begin{solution}
\(\|f_n\|_\infty\le1/n\to0\), but \(f_n'(x)=\cos(nx)\) has no pointwise limit at many points and no uniform limit.
\end{solution}

\begin{exercise}[Pointwise series manipulation]\label{exr:ii12-expand-12}
Why is pointwise convergence of a function series insufficient for termwise integration in general?
\end{exercise}
\begin{hint}
Partial sums may concentrate error on moving sets.
\end{hint}
\begin{solution}
Without uniform or another dominating control, integrals of the tails need not tend to zero even if tails vanish pointwise.
\end{solution}

\subsection*{Applications and investigations}

\begin{exercise}[Certified integral approximation]\label{exr:ii12-expand-13}
If a numerical surrogate \(\tilde f\) has uniform error \(\eta\) on an interval of length \(L\), what integral certificate follows?
\end{exercise}
\begin{hint}
Integrate the absolute error.
\end{hint}
\begin{solution}
The integral error is at most \(L\eta\).
\end{solution}

\begin{exercise}[Parameter differentiation]\label{exr:ii12-expand-14}
What kind of control should you seek before differentiating a limit of approximating functions?
\end{exercise}
\begin{hint}
Pointwise convergence is too weak; control derivatives uniformly and anchor function values.
\end{hint}
\begin{solution}
A standard sufficient pattern is uniform convergence of derivatives plus convergence at one base point.
\end{solution}

\subsection*{Challenge problems}

\begin{exercise}[Series derivative theorem]\label{exr:ii12-expand-15}
State a practical sufficient condition for differentiating \(\sum f_n\) termwise on \([a,b]\).
\end{exercise}
\begin{hint}
Control the derivative series uniformly and make the original series converge at one point.
\end{hint}
\begin{solution}
If \(\sum f_n'(x)\) converges uniformly and \(\sum f_n(x_0)\) converges at one point, then the original series converges uniformly to a differentiable sum with derivative \(\sum f_n'\).
\end{solution}

\begin{exercise}[Interchange as operator continuity]\label{exr:ii12-expand-16}
Explain why integration is continuous in the sup norm but differentiation is not.
\end{exercise}
\begin{hint}
Compare the bounds for the two operators.
\end{hint}
\begin{solution}
Integration satisfies \(|\int g|\le(b-a)\|g\|_\infty\). No bound of \(\|f'\|_\infty\) by a constant times \(\|f\|_\infty\) holds on \(C^1\); small-amplitude high-frequency functions show differentiation can amplify oscillation.
\end{solution}
% END VOL02-EXPANSION II12-exercises-01
\section*{Chapter summary}
The chapter's tools now form part of the canonical Volume II analysis toolkit and will be used in later approximation, fixed-point, and convergence arguments.
